% Nouveau EGA V (cote 25): archive cote 25, site page 151 (PDF page 152).
% English translation of the checked French annotated master.
% AI-translated by GPT-6 Astra (OpenAI, ChatGPT Pro mode).
\scanpage{25/151}
% Both long diagonal marks have unresolved function; both text blocks remain live.
\mspage{\hand{12}}
\var{\margin{\unsure{$\diagup$}}}{\placedmargin{\unsure{$\diagup$}}}{marginal note placed}{ega25_25_151-v0001}

projective of rank m, $u{:}E \hand{\to} I$ a homomorphism, giving a homomorphism
\[ u\otimes id_{A/I} : E\otimes A/I \hand{\longrightarrow} I/I^2 \ . \]
\var{\xout{Conditions Put}}{}{typewriter strike-out}{ega25_25_151-v0002} \var{\inserted[typed]{Let}}{Let}{insertion}{ega25_25_151-v0003} $B = A/u(E)A$, equipped with the augmentation $B\hand{\to} B/J = A/I$
(where $J = I/u(E)A$). The following two conditions are equivalent:

\var{\margin{$\|$}}{\placedmargin{$\|$}}{marginal note placed}{ega25_25_151-v0004} (i) The \var{\xout{augmented}}{}{typewriter strike-out}{ega25_25_151-v0005} ring B \var{\xout{B/I is}}{}{typewriter strike-out}{ega25_25_151-v0006} has an augmentation $B\hand{\to} B/J$ \var{\inserted{of type}}{of type}{insertion}{ega25_25_151-v0007} ordinary
quadratic, of rank $n-(m-1)= n-m+1$ (the \hand{``}\uuline{rank of the augmentation}\hand{''}
is \var{\inserted{by definition}}{by definition}{insertion}{ega25_25_151-v0008} that of the projective Module $J/J^2$).

\var{\margin{$|$}}{\placedmargin{$|$}}{marginal note placed}{ega25_25_151-v0009} (ii) $K = \mathrm{Ker}\ u\otimes id_{A/I}$ is an invertible Module, and the
\uline{natural} homomor\uline{phism} $K \hand{\longrightarrow} \mathrm{Sym}^2(J/J^2) = \mathrm{Sym}^2( (I/I^2)/\mathrm{Im}\ E\otimes A/I )$ is
``nondegenerate''.

It might be clearer to introduce a notation $V = I/I^2$,
and $W = J/J^2 = V / \mathrm{Im}\ E\otimes A/I$. The homomorphism described as natural
in (ii) is obtained by observing that K maps in any case into
\var{\margin{$\|$}}{\placedmargin{$\|$}}{marginal note placed}{ega25_25_151-v0010} $I^2/(I^3 +$ \var{\xout{Ixm(E)}}{}{typewriter strike-out}{ega25_25_151-v0011} $u(IE)) = (I^2/I^3)/\mathrm{Im}\ IE \hand{\simeq} \mathrm{Sym}^2(I/I^2)/\ \mathrm{Im}\ (I/I^2)\otimes(E\otimes A/$
% The source formula is clipped at the right edge; nothing is supplied.
where $I/I^2 \otimes(E \otimes A/I) \hand{\to} \mathrm{Sym}^2(I/I^2)$ is defined via the canonical
multiplication homomorphism $I/I^2 \otimes I/I^2 \hand{\to} \mathrm{Sym}^2(I/I^2 )$. With these descriptions
made explicit, the proof of Prop. 9 is almost trivial.

\var{\margin{\unsure{$\diagup$}}}{\placedmargin{\unsure{$\diagup$}}}{marginal note placed}{ega25_25_151-v0012}
\uline{Corollary} 10. Let A be a regular local ring of dimension n, $\underline{m}$ its
maximal ideal, \var{\inserted[typed]{$k=A/\underline{m}$,}}{$k=A/\underline{m}$,}{insertion}{ega25_25_151-v0013} E a free A-module of rank m, $u{:}E \hand{\to} \underline{M}$ a homomorphism,
\var{\xout{Equivalent conditions:}}{}{typewriter strike-out}{ega25_25_151-v0014} $B = A/u(E)A$, $\underline{n} = \underline{m}/u(E)A$ the maximal ideal of B,
$V = \underline{m}/\underline{m}^2$, $W = \underline{n}/\underline{n}^2 =$ \var{\xout{km/}}{}{typewriter strike-out}{ega25_25_151-v0015} $V/\ \mathrm{Im}\ E\otimes k$. Equivalent conditions:

\var{\margin{?\}}}{\placedmargin{?\}}}{marginal note placed}{ega25_25_151-v0016} (i) B has dimension n-m (i.e., in terms of a basis \var{\xout{of}}{}{typewriter strike-out}{ega25_25_151-v0017} $(\hand{\phi}_i)$ of
\var{\margin{$\|$}}{\placedmargin{$\|$}}{marginal note placed}{ega25_25_151-v0018} E, the images of the $\hand{\phi}_i$ in A form an \tikz[baseline=(t.base)]{\node[draw,ellipse,handwritten,text=black,inner sep=1pt] (t) {A-régulière};} sequence and \var{\xout{is}}{}{typewriter strike-out}{ega25_25_151-v0020} the
\var{\xout{singularity qu}}{}{typewriter strike-out}{ega25_25_151-v0021} closed point of Spec(B) is an \uline{ordinary quadratic singularity}
(by which we mean that the augmentation $B\hand{\to} B/\underline{n}$ is \var{\inserted{of type}}{of type}{insertion}{ega25_25_151-v0022} ordinary
quadratic).

(ii) There exists a basis $\hand{\phi}_1 ,\ldots., \hand{\phi}_m$ of E such that \var{\inserted{a)}}{a)}{insertion}{ega25_25_151-v0023}$A/$\hand{$\sum_{i=1}^{m-1}$}$u(\hand{\phi} i)A$

\mspage{\hand{151}}
\tnote{25-151-1}{The clipped formula ending in A/ is left incomplete. The underlined suffix of homomorphism and the adjective natural correspond to the underlined French phisme naturel. The circled label A-régulière is reproduced literally and means A-regular. The canceled t after les has no English effect. The canceled fragment qu remains unexpanded after singularity; no completion of that French word is supplied. The capital underlined M in the target of the Corollary's homomorphism and the unmatched parenthesis in (i) remain as in the source.}
